% ==============================================================================
% ABSTRACT (inglés) — Página preliminar de acceso directo (APA 7 Nivel 1)
% Fuente: texto_formatter/abstract.txt (formateado en APA 7)
% ==============================================================================

\section*{Abstract}
\addcontentsline{toc}{section}{Abstract}

This project aims to develop a comprehensive information system to optimize the application, reception, and selection process for candidates of the \enquote{Beca Comedor} (Dining Hall Scholarship) at the Universidad Autónoma Gabriel René Moreno (UAGRM). Currently, the management of this benefit presents deficiencies because the evaluation of applications and profiles is carried out strictly manually. This situation generates bottlenecks, prolonged response times, and a significant operational workload for the staff of the University Directorate of Social Welfare and Health (DUBSS). Additionally, the student population lacks a unified digital channel that facilitates remote applications and transparent access to the daily menu.

To solve this problem, the development of a technological platform structured into two components is proposed: a web portal oriented towards administrative management and a mobile application for students. The web portal will allow the DUBSS to centralize information, manage profiles, and administer the food service; meanwhile, the mobile application will facilitate university students in applying for the scholarship and checking the updated dining hall menu.

The core of the solution lies in the integration of Artificial Intelligence (AI) for automated pre-selection. This model will process the applicants' data, efficiently discarding those applications that do not meet the minimum requirements, massively reducing the burden of human review. Complementarily, a Business Intelligence (BI) module will be incorporated to provide interactive dashboards to the authorities. These analytical tools will allow the monitoring of key indicators regarding the efficiency of the pre-selection and the use of the dining hall, facilitating informed decision-making. Altogether, the project modernizes and streamlines the administration of student welfare.

\indent\textit{Keywords: Information System, University Dining Hall, Artificial Intelligence, Business Intelligence, Process Automation, Mobile Application, Administrative Management.}